\section{结论}
\label{sec:conclusion}

本文提出了 PlanDiffusion，一种以图结构为中间表示的文本驱动住宅平面图生成框架。
通过将生成任务分解为自回归图拓扑生成（Stage~1）和节点坐标扩散（Stage~2）两个阶段，
本方法在保证房间连通拓扑合法性的同时，实现了从自然语言到矢量平面图的端到端生成。

为支撑模型训练，我们构建了 PlanText-150K 数据集，
包含 149,609 条具有高质量自然语言描述的平面图样本，
是目前规模最大的带文本标注住宅平面图数据集之一。

\todo{实验结果完成后补充定量结论和局限性分析。}

未来工作将探索更长文本条件下的生成质量提升、
多楼层平面图生成，以及基于用户反馈的交互式编辑能力。
